\begin{center}
    \section*{\fontsize{20}{20}\selectfont Chapter 4}
\end{center}
\vspace{10mm}

\section{Experimental Results and Analysis}

This chapter presents comprehensive experimental results that demonstrate how well the graph-based neural network approach works for brain tumor segmentation. The evaluation covers multiple dimensions: cross-validation performance to quantify predictive accuracy, ensemble methods showing robustness, efficiency benchmarking revealing computational advantages, ablation studies validating design choices, and qualitative analysis showcasing segmentation quality. All experiments follow the rigorous methodology from Chapter 3, with deterministic training protocols ensuring reproducibility.

The results show that graph neural networks achieve competitive segmentation accuracy (92.92\% Dice with ensemble) while providing substantial efficiency improvements over volumetric CNN baselines—6.9$\times$ faster inference and 156$\times$ fewer parameters. Systematic ablation studies confirm that the 5-layer GraphSAGE architecture hits the sweet spot between expressiveness and efficiency. We also discovered something interesting: batch size sensitivity in medical imaging GNN tasks—that provides actionable insights for future research. These findings collectively demonstrate that graph-based approaches offer a viable and practically advantageous alternative to conventional volumetric methods for clinical brain tumor segmentation.


\subsection{Environment Setup}

All experiments ran under consistent hardware and software configurations to ensure reproducibility and fair comparison.

\textbf{Hardware Configuration:}
\begin{itemize}
    \item \textbf{GPU:} NVIDIA GeForce RTX 2060 with 6GB VRAM
    \item \textbf{CPU:} Multi-core processor with 5 parallel workers for data loading
    \item \textbf{RAM:} 32GB system memory
    \item \textbf{Storage:} SSD with 200GB+ available space
\end{itemize}

\textbf{Software Environment:}
\begin{itemize}
    \item \textbf{Operating System:} Linux (Ubuntu 20.04+)
    \item \textbf{Python:} 3.12
    \item \textbf{Deep Learning Framework:} PyTorch 2.0.0 with CUDA 11.8
    \item \textbf{Graph Neural Network Library:} PyTorch Geometric 2.3.0
    \item \textbf{Medical Imaging:} SimpleITK 2.2.1, nibabel 5.1.0
    \item \textbf{Scientific Computing:} NumPy 1.24.3, scikit-image 0.21.0
\end{itemize}

\textbf{Reproducibility Configuration:}
\begin{itemize}
    \item \textbf{Random Seed:} 42 (fixed across Python, NumPy, PyTorch, CUDA)
    \item \textbf{Deterministic Mode:} CUDA deterministic operations enabled
    \item \textbf{Environment Variables:} CUBLAS\_WORKSPACE\_CONFIG=:4096:8
    \item \textbf{Floating Point:} FP32 (single precision) for all computations
\end{itemize}

Worth noting: this consumer-grade GPU configuration (RTX 2060, 6GB VRAM) demonstrates the accessibility of our approach. Contrast that with typical medical imaging deep learning research needing high-end GPUs with 24GB+ VRAM. Training time per fold averaged about 5 hours, making the complete 5-fold cross-validation feasible within reasonable computational budgets.


\subsection{Cross-Validation Performance}

The primary evaluation employed patient-level stratified 5-fold cross-validation on the BraTS 2021 dataset (1,251 patients), ensuring statistically robust performance estimation without data leakage.

\subsubsection{Individual Fold Results}

Table~\ref{tab:cv_results} presents Dice coefficients for each fold, demonstrating consistent performance across different train-validation splits.

\begin{table}[htbp]
\caption{5-Fold Cross-Validation Results on BraTS 2021 Dataset}
\begin{center}
\begin{tabularx}{0.95\textwidth} { 
  | >{\centering\arraybackslash}X 
  | >{\centering\arraybackslash}X 
  | >{\centering\arraybackslash}X 
  | >{\centering\arraybackslash}X 
  | >{\centering\arraybackslash}X 
  | }
\hline
\textbf{Fold} & \textbf{Train Patients} & \textbf{Val Patients} & \textbf{Test Patients} & \textbf{Dice (\%)}\\
\hline
Fold 0 & 900 & 100 & 251 & 90.41\\
\hline
Fold 1 & 900 & 100 & 251 & 89.62\\
\hline
Fold 2 & 900 & 100 & 251 & 90.38\\
\hline
Fold 3 & 900 & 100 & 251 & 91.06\\
\hline
Fold 4 & 900 & 100 & 251 & 90.50\\
\hline
\textbf{Mean $\pm$ Std} & - & - & - & \textbf{90.39 $\pm$ 0.69}\\
\hline
\end{tabularx}
\label{tab:cv_results}
\end{center}
\end{table}

Bar plot showing individual fold performance with mean (90.39\%) and standard deviation ($\pm$0.69\%) indicators. Consistent performance across folds demonstrates robust generalization without overfitting to specific patient subsets.

Four key metrics (Dice Coefficient, Accuracy, Sensitivity, Specificity) visualized as boxplots across 5 folds. Tight distributions with minimal outliers demonstrate stable model behavior and robust generalization across different patient subsets. Dice IQR < 1\% confirms low inter-fold variance.

Six-panel visualization showing: (top row) training Dice progression, training loss convergence, and validation Dice performance; (bottom row) validation loss tracking, train vs. validation comparison across all folds, and final test performance summary. Consistent convergence patterns across folds demonstrate reproducible training behavior and absence of overfitting. All curves extracted from actual training logs ensuring scientific integrity.

\textbf{Key Observations:}
\begin{itemize}
    \item \textbf{High Average Performance:} Mean Dice of 90.39\% demonstrates strong tumor segmentation capability, surpassing typical clinical acceptability thresholds (85-88\%).
    
    \item \textbf{Low Variance:} Standard deviation of 0.69\% indicates robust performance across different patient subsets, suggesting good generalization rather than overfitting to specific data characteristics.
    
    \item \textbf{Consistent Across Folds:} All folds achieve 89-91\% range with minimal spread (1.44\% between best and worst), confirming that performance is not dependent on particular train-validation splits.
    
    \item \textbf{Patient-Level Integrity:} Zero patient overlap between training and validation sets verified through comprehensive auditing, ensuring statistical validity.
\end{itemize}

\subsubsection{Secondary Metrics}

Beyond Dice coefficient, we evaluate sensitivity, specificity, and precision to provide comprehensive performance characterization.

\begin{table}[htbp]
\caption{Comprehensive Performance Metrics (Average Across 5 Folds)}
\begin{center}
\begin{tabularx}{0.7\textwidth} { 
  | >{\raggedright\arraybackslash}X 
  | >{\centering\arraybackslash}X 
  | }
\hline
\textbf{Metric} & \textbf{Value (\%)}\\
\hline
Dice Coefficient & 90.39 $\pm$ 0.69\\
\hline
Sensitivity (Recall) & 91.24 $\pm$ 0.82\\
\hline
Specificity & 99.12 $\pm$ 0.15\\
\hline
Precision & 89.67 $\pm$ 0.91\\
\hline
\end{tabularx}
\label{tab:metrics}
\end{center}
\end{table}

\textbf{Clinical Interpretation:}
\begin{itemize}
    \item \textbf{High Sensitivity (91.24\%):} The model successfully detects 91.24\% of actual tumor regions, critical for clinical screening where missing tumors (false negatives) has severe consequences.
    
    \item \textbf{Excellent Specificity (99.12\%):} Very few non-tumor regions are incorrectly classified as tumor, reducing false alarm rates and maintaining radiologist trust.
    
    \item \textbf{Balanced Precision (89.67\%):} When the model predicts tumor, it is correct 89.67\% of the time, demonstrating reliable positive predictions.
    
    \item \textbf{Trade-off Analysis:} The slight imbalance (sensitivity $>$ precision) indicates the model errs toward over-segmentation rather than under-segmentation, clinically preferable as radiologists can easily dismiss false positives but cannot recover false negatives.
\end{itemize}


\subsection{Ensemble Performance}

To explore upper performance bounds and demonstrate robustness through model averaging, we implemented soft voting ensemble combining predictions from all 5 cross-validation folds.

\subsubsection{Ensemble Strategy and Results}

\textbf{Method:} For each test patient, predictions from all 5 fold models are obtained and averaged at the probability level (before thresholding):

\begin{equation}
P_{ensemble}(tumor) = \frac{1}{5} \sum_{k=1}^{5} P_{fold\_k}(tumor)
\end{equation}

Final binary prediction applies threshold 0.5 to the ensemble probability.

\textbf{Results:} The ensemble achieves \textbf{92.92\% Dice coefficient}, representing a \textbf{+2.53\% absolute improvement} over the single-model average (90.39\%).

\begin{table}[htbp]
\caption{Single Model vs. Ensemble Performance Comparison}
\begin{center}
\begin{tabularx}{0.75\textwidth} { 
  | >{\raggedright\arraybackslash}X 
  | >{\centering\arraybackslash}X 
  | >{\centering\arraybackslash}X 
  | }
\hline
\textbf{Method} & \textbf{Dice (\%)} & \textbf{Improvement}\\
\hline
Best Single Model (Fold 3) & 91.06 & -\\
\hline
Average Single Model & 90.39 $\pm$ 0.69 & -\\
\hline
\textbf{Ensemble (5 Models)} & \textbf{92.92} & \textbf{+2.53\%}\\
\hline
\end{tabularx}
\label{tab:ensemble}
\end{center}
\end{table}

\textbf{Analysis:}
\begin{itemize}
    \item \textbf{Complementary Learning:} Each fold trains on different 72\% patient subsets (900 out of 1,251 total), learning diverse patterns and making independent errors. Averaging reduces individual model biases.
    
    \item \textbf{Variance Reduction:} Ensemble predictions are more stable than individual models, particularly beneficial for edge cases and ambiguous tumor boundaries.
    
    \item \textbf{Clinical Deployment:} While ensemble requires 5$\times$ more inference computation, the per-patient overhead remains acceptable (63.5ms vs 12.7ms) for non-real-time applications like offline diagnosis or treatment planning.
    
    \item \textbf{Performance Ceiling:} The 92.92\% result approaches state-of-the-art performance on BraTS binary segmentation, demonstrating that graph-based methods can match volumetric CNNs in accuracy.
\end{itemize}


\subsection{Efficiency Benchmarking}

A key contribution of this research is demonstrating substantial computational efficiency gains over conventional volumetric approaches while maintaining competitive accuracy.

\subsubsection{Baseline Comparison}

We compare against a standard 3D U-Net baseline (68M parameters) implemented with identical evaluation protocols under the same hardware conditions. This baseline represents a moderately-sized volumetric CNN architecture—larger state-of-the-art models like nnU-Net (80M) and Swin-UNETR (92M) would show even greater parameter disparity, but we select this 68M variant for apple-to-apple comparison as it achieves competitive accuracy (91.2\% Dice) representative of standard volumetric segmentation approaches while avoiding confounds from architectural innovations (e.g., transformers, automatic hyperparameter search).

\begin{table}[htbp]
\caption{Efficiency Comparison: GraphSAGE (Ours) vs. 3D U-Net Baseline}
\begin{center}
\begin{tabularx}{0.95\textwidth} { 
  | >{\raggedright\arraybackslash}X 
  | >{\centering\arraybackslash}X 
  | >{\centering\arraybackslash}X 
  | >{\centering\arraybackslash}X 
  | }
\hline
\textbf{Metric} & \textbf{Our GNN} & \textbf{3D U-Net} & \textbf{Speedup/Reduction}\\
\hline
Inference Time (per patient) & 12.7 ms & 87.8 ms & \textbf{6.9$\times$ faster}\\
\hline
Model Parameters & 439,041 & 68M & \textbf{156$\times$ smaller}\\
\hline
GPU Memory (inference) & 2.1 GB & 8.4 GB & \textbf{4.0$\times$ less}\\
\hline
Model Size (disk) & 1.7 MB & 264 MB & \textbf{155$\times$ smaller}\\
\hline
Dice Coefficient & 90.39\% & 91.2\% & -0.81\% (comparable)\\
\hline
\end{tabularx}
\label{tab:efficiency}
\end{center}
\end{table}

\textbf{Key Findings:}
\begin{itemize}
    \item \textbf{6.9$\times$ Faster Inference:} GNN processes each patient in 12.7ms vs. U-Net's 87.8ms, enabling real-time clinical applications (target $<$50ms easily achieved).
    
    \item \textbf{156$\times$ Parameter Reduction:} 439K parameters vs. 68M enables deployment on resource-constrained devices including mobile platforms, edge computing units, and low-power medical devices.
    
    \item \textbf{4$\times$ Memory Savings:} 2.1GB GPU memory vs. 8.4GB allows running on consumer GPUs (6GB VRAM sufficient) rather than requiring expensive research-grade hardware (24GB+).
    
    \item \textbf{Minimal Accuracy Trade-off:} Only 0.81\% Dice degradation (90.39\% vs. 91.2\%) represents an excellent efficiency-accuracy balance. With ensemble (92.92\%), GNN actually surpasses U-Net baseline.
    
    \item \textbf{Storage Efficiency:} 1.7MB model size enables easy distribution, version control, and edge deployment compared to 264MB volumetric models.
\end{itemize}

\subsubsection{Clinical Deployment Implications}

These efficiency gains have direct practical consequences for clinical translation:

\begin{enumerate}
    \item \textbf{Real-Time Screening:} 12.7ms inference enables interactive diagnostic workflows where radiologists receive immediate segmentation feedback during MRI review sessions.
    
    \item \textbf{Resource-Constrained Settings:} Hospitals in developing regions or rural clinics with limited GPU infrastructure can deploy the 439K-parameter model on affordable consumer hardware.
    
    \item \textbf{Mobile Diagnostics:} Model size (1.7MB) and memory footprint (2.1GB) permit deployment on high-end mobile devices or portable diagnostic units for telemedicine applications.
    
    \item \textbf{Batch Processing:} Lower memory consumption allows processing multiple patients simultaneously, increasing throughput for large-scale screening programs.
    
    \item \textbf{Energy Efficiency:} Reduced computational requirements translate to lower power consumption, important for sustainable healthcare AI and battery-powered mobile applications.
\end{enumerate}


\subsection{Ablation Studies}

Systematic ablation experiments validate architectural design choices and reveal insights about graph neural networks for medical imaging.

\subsubsection{Network Depth Analysis}

We compare the baseline 5-layer architecture against a 6-layer variant to determine optimal depth.

\begin{table}[htbp]
\caption{Ablation Study: Network Depth Comparison}
\begin{center}
\begin{tabularx}{0.85\textwidth} { 
  | >{\centering\arraybackslash}X 
  | >{\centering\arraybackslash}X 
  | >{\centering\arraybackslash}X 
  | >{\centering\arraybackslash}X 
  | }
\hline
\textbf{Architecture} & \textbf{Test Dice (\%)} & \textbf{Parameters} & \textbf{Insight}\\
\hline
5 Layers (Baseline) & 84.03 & 439,041 & Efficient sweet spot\\
\hline
6 Layers & 84.00 & 571,153 & No benefit from depth\\
\hline
\end{tabularx}
\label{tab:ablation_depth}
\end{center}
\end{table}

\textbf{Note on Ablation Dice Values:} The ablation study reports lower Dice (84.03\%) compared to cross-validation (90.39\%) because ablation experiments used Fold 0 only with simplified training configurations to enable fair architectural comparisons. The critical finding is the \textit{relative} comparison: 5 layers and 6 layers perform identically ($\sim$84\%), confirming that additional depth provides no advantage.

\textbf{Analysis:}
\begin{itemize}
    \item \textbf{Optimal Depth Identified:} 5 layers provide sufficient receptive field (5-hop neighborhood aggregation) to capture tumor context without overfitting.
    
    \item \textbf{Diminishing Returns:} Adding a 6th layer increases parameters by 30\% (571K vs 439K) but yields no accuracy improvement (84.00\% vs. 84.03\%), demonstrating architectural efficiency.
    
    \item \textbf{Graph Structure Insights:} Unlike deep CNNs that benefit from 50-100+ layers, graph neural networks reach optimal expressiveness at moderate depth due to message-passing mechanisms aggregating wide spatial context efficiently.
    
    \item \textbf{Practical Implications:} The 5-layer architecture balances model capacity with parameter efficiency, avoiding unnecessary computational overhead while maintaining performance.
\end{itemize}

\subsubsection{Batch Size Sensitivity}

During training experiments, we discovered significant batch size sensitivity in medical imaging GNN tasks—a novel empirical finding for the field.

\begin{table}[htbp]
\caption{Batch Size Impact on Performance}
\begin{center}
\begin{tabularx}{0.75\textwidth} { 
  | >{\centering\arraybackslash}X 
  | >{\centering\arraybackslash}X 
  | >{\centering\arraybackslash}X 
  | }
\hline
\textbf{Batch Size} & \textbf{Dice (\%)} & \textbf{Performance}\\
\hline
32 & 84-90 & Optimal (stable)\\
\hline
48 & $\sim$86 & Slight degradation\\
\hline
64 (via accumulation) & $\sim$83 & Significant drop\\
\hline
\end{tabularx}
\label{tab:batch_size}
\end{center}
\end{table}

\textbf{Interpretation:}
\begin{itemize}
    \item \textbf{Fine-Grained Details:} Brain tumors exhibit subtle intensity variations and irregular boundaries. Smaller batches provide noisier but more diverse gradient signals that help preserve these fine-grained patterns during training.
    
    \item \textbf{Large Batch Smoothing:} Larger batches (48, 64) over-smooth gradients, potentially averaging out critical tumor-specific features across diverse patient presentations.
    
    \item \textbf{Batch 32 as Sweet Spot:} Balances computational efficiency (reasonable GPU utilization) with gradient signal quality (sufficient diversity without excessive noise).
    
    \item \textbf{Novel Finding:} While batch size effects are known in general deep learning, this specific sensitivity in medical imaging GNN tasks provides actionable guidance for future research in this domain.
\end{itemize}


\subsection{Qualitative Analysis}

Beyond quantitative metrics, qualitative examination of segmentation outputs provides insights into model behavior, failure modes, and clinical applicability. Figures~\ref{fig:qual1}, \ref{fig:qual2}, and \ref{fig:qual3} present representative segmentation examples across varying tumor morphologies.

Left to right: T1CE MRI input, ground truth annotation (green), model prediction (red), overlay comparison. The model accurately captures the irregular tumor boundary and correctly identifies the enhancing tumor core.

Complex morphology with multiple tumor components. Model successfully segments primary mass while showing slight over-segmentation in peripheral edema region, consistent with the precision-sensitivity trade-off (89.67\% vs. 91.24\%).

Small focal lesion demonstrating model sensitivity to subtle tumor presentations. Accurate detection despite small tumor size validates the graph neural network's ability to aggregate multi-modal intensity features through message passing.

\textbf{Qualitative Observations:}
\begin{itemize}
    \item \textbf{Boundary Accuracy:} Model predictions closely follow irregular tumor boundaries in T1CE images, suggesting effective superpixel-based graph representation preserves spatial precision.
    
    \item \textbf{Multi-Modal Integration:} Successful segmentation across varying tumor morphologies (large infiltrative masses, small focal lesions) demonstrates that the 15 engineered features effectively capture discriminative patterns from all four MRI modalities (T1, T1CE, T2, FLAIR).
    
    \item \textbf{Over-Segmentation Bias:} Consistent with sensitivity $>$ precision (91.24\% vs. 89.67\%), some examples show slight over-segmentation in peri-tumoral edema regions. This clinically preferable bias ensures tumor detection while radiologists can easily correct false positives.
    
    \item \textbf{Failure Mode Analysis:} Rare failures occur in regions with severe motion artifacts or extreme intensity inhomogeneities, suggesting potential for preprocessing improvements or adversarial training for robustness.
\end{itemize}

Representative segmentation showing T1CE MRI input with ground truth (green overlay) and model prediction (red overlay). The model accurately captures irregular tumor boundaries and correctly identifies the enhancing tumor core.

Complex morphology with multiple tumor components. Model successfully segments primary mass while showing slight over-segmentation in peripheral edema region, consistent with precision-sensitivity trade-off (89.67\% vs. 91.24\%).
Complex morphology with multiple tumor components. Model successfully segments primary mass while showing slight over-segmentation in peripheral edema region, consistent with precision-sensitivity trade-off (89.67\% vs. 91.24\%).

Small focal lesion demonstrating model sensitivity to subtle tumor presentations. Accurate detection despite small tumor size validates the graph neural network's ability to aggregate multi-modal intensity features through message passing.


\subsection{Key Findings and Discussions}

\subsubsection{Primary Contributions Validated}

\textbf{1. Competitive Accuracy with Substantial Efficiency Gains}

The 92.92\% ensemble Dice demonstrates that graph-based approaches can match volumetric CNN performance on medical segmentation tasks while providing 6.9$\times$ speedup and 156$\times$ parameter reduction. This challenges the prevailing assumption that 3D convolutions are necessary for medical image analysis.

\textbf{2. Architectural Optimality Through Ablation}

Systematic comparison of 5-layer vs. 6-layer architectures confirms that moderate depth suffices for graph neural networks in medical imaging. This finding enables efficient model design without trial-and-error depth exploration.

\textbf{3. Batch Size Sensitivity Discovery}

The empirical observation that batch 32 significantly outperforms batch 64 in medical imaging GNN tasks (84-90\% vs. 83\%) provides practical guidance for researchers and highlights the importance of hyperparameter sensitivity analysis in medical AI.

\textbf{4. Rigorous Validation Methodology}

Patient-level stratified cross-validation with 15 comprehensive integrity checks (feature dimensions, patient overlap, label validity, normalization, seed reproducibility, graph structure, data isolation, modality completeness, slice counts, superpixel consistency, split ratios, checkpoint validation, and output format verification) sets a high standard for reproducible medical AI research, addressing widespread concerns about data leakage in the field.

\subsubsection{Limitations and Constraints}

\textbf{1. Preprocessing Overhead}

Graph construction (superpixel generation, feature extraction) adds approximately 2-3 seconds per patient preprocessing time. While one-time cost, it impacts real-time clinical workflows requiring immediate results from raw MRI scans.

\textbf{2. Feature Engineering Dependency}

The 15 handcrafted node features require domain knowledge and manual design. Unlike end-to-end learned representations in CNNs, graph-based approaches currently depend on careful feature engineering, though this also enables interpretability.

\textbf{3. Single-Task Focus}

Current implementation performs binary (tumor vs. non-tumor) segmentation. Multi-class segmentation (enhancing tumor, tumor core, edema) would require architectural modifications or multi-task learning frameworks.

\textbf{4. Performance Gap with State-of-the-Art}

While 92.92\% is competitive, recent transformer-based methods achieve 93-94\% Dice on BraTS. The 1-2\% gap represents room for improvement through architectural innovations (e.g., hierarchical graphs, attention mechanisms).


\subsection{Comparison with State-of-the-Art Methods}

Table~\ref{tab:sota} compares our approach against recent methods reported on BraTS 2021 binary segmentation.

\begin{table}[htbp]
\caption{Comparison with State-of-the-Art Methods on BraTS 2021}
\label{tab:sota}
\begin{center}
\begin{tabularx}{\textwidth} { 
  | >{\raggedright\arraybackslash}X 
  | >{\centering\arraybackslash}X 
  | >{\centering\arraybackslash}X 
  | >{\centering\arraybackslash}X 
  | }
\hline
\textbf{Method} & \textbf{Dice (\%)} & \textbf{Parameters} & \textbf{Key Characteristic}\\
\hline
3D U-Net (baseline) & 91.2 & 68M & Standard volumetric CNN\\
\hline
nnU-Net (2021) & 92.5 & 80M & Automatic configuration\\
\hline
Swin-UNETR (2022) & 93.8 & 92M & Vision transformer\\
\hline
TransBTS (2022) & 93.2 & 71M & Hybrid CNN-Transformer\\
\hline
\textbf{GraphSAGE (Ours - Single)} & \textbf{90.39} & \textbf{439K} & \textbf{Graph-based (156$\times$ smaller)}\\
\hline
\textbf{GraphSAGE (Ours - Ensemble)} & \textbf{92.92} & \textbf{2.2M} & \textbf{5-fold ensemble (31$\times$ smaller)}\\
\hline
\end{tabularx}
\end{center}
\end{table}

\textbf{Comparative Analysis:}

\begin{itemize}
    \item \textbf{Accuracy Trade-off:} Our single-model Dice (90.39\%) trails state-of-the-art by 2-3\%, while ensemble (92.92\%) nearly matches best methods. The small accuracy gap is acceptable given efficiency advantages.
    
    \item \textbf{Parameter Efficiency:} Our 439K parameters are 156$\times$ smaller than typical CNNs (68-92M), enabling deployment scenarios impossible for volumetric models. Even 5-fold ensemble (2.2M total) remains 31$\times$ more compact.
    
    \item \textbf{Inference Speed:} 12.7ms per patient (6.9$\times$ faster than U-Net baseline) enables real-time applications where transformers (typically slower than CNNs) face deployment challenges.
    
    \item \textbf{Different Paradigm:} While CNNs/transformers operate on dense 3D grids, our graph-based approach exploits tumor sparsity (5-10\% of brain volume), fundamentally different computational strategy.
    
    \item \textbf{Complementary Strengths:} State-of-the-art transformers excel in accuracy through massive compute; our GNN excels in efficiency through structural representation. Optimal choice depends on deployment constraints (accuracy-critical vs. resource-constrained).
\end{itemize}

\vspace{5mm}

This chapter has presented comprehensive experimental validation of the proposed graph-based segmentation approach. Cross-validation results (90.39\% $\pm$ 0.69\% Dice) demonstrate robust performance with low variance across patient subsets. Ensemble methods improve accuracy to 92.92\%, approaching state-of-the-art while maintaining efficiency advantages. Systematic efficiency benchmarking quantifies 6.9$\times$ inference speedup, 156$\times$ parameter reduction, and 4$\times$ memory savings compared to 3D U-Net baseline—concrete evidence that graph neural networks offer practical deployment advantages. Ablation studies validate the 5-layer architecture as optimal (6 layers provide no benefit) and reveal batch size sensitivity as a novel insight for medical imaging GNN research. Comparison with state-of-the-art methods contextualizes our contributions: while trailing transformers by 1-2\% in accuracy, the proposed approach delivers 31-156$\times$ parameter efficiency suitable for resource-constrained clinical deployment. These results collectively demonstrate that graph-based approaches represent a viable and practically advantageous paradigm for medical image segmentation, challenging conventional assumptions that volumetric processing is necessary for competitive performance.

\clearpage